\documentclass[11pt,letterpaper]{article}

\usepackage[margin=0.72in]{geometry}
\usepackage[T1]{fontenc}
\usepackage{lmodern}
\usepackage[hidelinks]{hyperref}
\usepackage{microtype}
\setlength{\parindent}{0pt}
\setlength{\parskip}{8pt}
\pagestyle{empty}

\begin{document}

\textbf{\Large Piter Garcia}\\
Rochester, NY \enspace|\enspace Available for NYC relocation beginning January 2027\\
\href{mailto:pgarcia8@u.rochester.edu}{pgarcia8@u.rochester.edu} \enspace|\enspace
+1 (646) 288-3300\\
\href{https://linkedin.com/in/garciapiterz}{linkedin.com/in/garciapiterz}
\enspace|\enspace
\href{https://garciapiterz.com}{garciapiterz.com}

\vspace{10pt}
August 19, 2026

Hiring Team\\
Explore Schools\\
Brooklyn, New York

\textbf{Re: Self Contained Elementary Special Education Teacher (SY 26--27)}

Dear Explore Schools Hiring Team:

I am writing in response to Brianna Bernard's outreach about the Self Contained
Elementary Special Education Teacher position. The opportunity stands out to me
because it brings together the areas I have deliberately built into my teacher
preparation: elementary instruction, disability and inclusion, culturally
responsive teaching, individualized support, and service to high-need school
communities.

I am completing an M.S. in Teaching Computer Science K--12 at the University of
Rochester as an NSF Robert Noyce Scholar, with graduation expected in December
2026. This spring, I completed K--5 student teaching at Pine Brook Elementary.
My placement moved from observation and individual or small-group support into
co-teaching and increasing responsibility for lesson design, instruction,
assessment, and classroom management. Across kindergarten through fifth grade,
I designed and adapted learning experiences using coding, robotics, digital
fluency, and hands-on STEM while adjusting pacing, grouping, visual supports,
task structure, and response options for varied learners.

My graduate preparation has included a substantial disability and inclusive-
education strand: Universal Design for Learning, differentiation, MTSS,
IEP-informed instruction, accommodations and modifications, assessment
adaptation, progress monitoring, behavior and self-regulation supports, and
collaboration among educators, related-service providers, and families. One of
the most important lessons from that work is that behavior and incomplete work
can be information about access, regulation, relationships, or instructional
design. At Pine Brook, I saw students who had been avoiding work become engaged,
successful participants when I changed the structure and entry point of the
task while preserving the learning goal. Some later began helping peers. That
experience reinforced my commitment to maintaining high expectations while
changing the supports and pathways students need to reach them.

I also bring an unusual technical and analytical foundation for an elementary
educator. Before entering teacher preparation, I worked in software, data, and
AI systems, and I am also completing an M.S. in Data Science at RIT. That
background has trained me to observe carefully, document patterns, validate
assumptions, and adjust from evidence without reducing a person to a metric. In
the classroom, I bring that same discipline to formative assessment, progress
monitoring, instructional iteration, and clear communication across a team.

Explore Schools' emphasis on rigorous learning, culturally responsive practice,
student belonging, and continuous educator growth is closely aligned with the
kind of teacher I am preparing to become. My direct classroom experience has
been in inclusive K--5 student teaching rather than as the teacher of record in
a self-contained 12:1:1 classroom, and I would represent that distinction
clearly. I would bring strong supervised preparation, a learner-centered and
data-informed approach, and a willingness to learn Explore Schools' specific
IEP, team, and classroom systems from experienced special educators and school
leaders.

I will remain based in Rochester while I finish my two master's degrees through
December 2026, and I can relocate to New York City and work fully onsite
beginning in January 2027. If Explore Schools is considering a mid-year start
for this opening, I would welcome the opportunity to discuss how my elementary,
inclusive-education, and technical background could support your students and
school community.

Thank you for your consideration.

Sincerely,\\[8pt]
\textbf{Piter Garcia}

\end{document}
